% ====================================
% CHƯƠNG 5: THIẾT KẾ CƠ SỞ DỮ LIỆU
% ====================================

\chapter{Thiết kế cơ sở dữ liệu}
\label{chap:thiet_ke_csdl}

\section{Mô hình thực thể - quan hệ (ERD)}
\label{sec:erd}

% Chèn sơ đồ ERD
% \begin{figure}[H]
%     \centering
%     \includegraphics[width=0.95\textwidth]{erd.png}
%     \caption{Sơ đồ thực thể -- quan hệ (ERD)}
%     \label{fig:erd}
% \end{figure}


\section{Mô hình quan hệ}
\label{sec:mo_hinh_quan_he}

% Chuyển đổi từ ERD sang mô hình quan hệ


\section{Danh sách bảng dữ liệu}
\label{sec:danh_sach_bang}

% Mô tả chi tiết từng bảng

\subsection{Bảng Users}
\label{subsec:bang_users}

\begin{table}[H]
    \centering
    \caption{Cấu trúc bảng Users}
    \label{tab:users}
    \begin{tabularx}{\textwidth}{|c|l|l|l|X|}
        \hline
        \textbf{STT} & \textbf{Tên cột} & \textbf{Kiểu dữ liệu} & \textbf{Ràng buộc} & \textbf{Mô tả} \\
        \hline
        1 & id & INT & PK, AI & Mã người dùng \\
        \hline
        2 & username & VARCHAR(50) & UNIQUE, NOT NULL & Tên đăng nhập \\
        \hline
        3 & email & VARCHAR(100) & UNIQUE, NOT NULL & Email \\
        \hline
        4 & password & VARCHAR(255) & NOT NULL & Mật khẩu (đã mã hóa) \\
        \hline
    \end{tabularx}
\end{table}

% Thêm các bảng khác tương tự


\section{Ràng buộc toàn vẹn}
\label{sec:rang_buoc}

% Mô tả các ràng buộc toàn vẹn dữ liệu
